% Chapter 05. Step 8 itself (trust weights, neutral imputation, one RL iteration with the referee) is
% in chapter 03 and is not repeated here (O17, option (a), decided 27 Nov 2026). This chapter holds
% decision D1: docs/thesis/notes/d1_scale_policy.md converted to LaTeX in D.2 without rewriting.
% Its references to the proposal ("Eq. (20)", "P14 (b)", ...) are links to chapter 03, which prints
% the proposal's numbers. The Markdown note stays in docs/thesis/notes/; change both together.
\chapter{Trust weighting}\label{ch:trust}

The trust weights and the neutral imputation with the trust-weighted mean are Step~8 of the
proposal, \cref{method:sec:trust}; they are not repeated here. This chapter holds decision D1 of the
implementation plan: which standard deviation scales the trust-weighted advantages. The proposal's
eq.~\eqref{method:eq:adv} uses the std of the trust-weighted returns (option A); option B, the std of
the raw returns, is an extra ablation (data-free roadmap, §0.1).

\section{D1: which std scales the trust-weighted advantages}\label{trust:sec:d1}

Companion to \code{src/rlinf\_modified/algorithms/trust.py} and \code{tests/test\_trust\_weighting.py}
(WP 1.2, 25 Nov 2026). Equation and proposition numbers refer to the proposal (27 Oct 2026).

\subsection{Setup}

Fix one group of $G$ candidates. Weights $w_{i,k}\in(0,1]$. Eq.~\eqref{method:eq:neutral}:
\[
m_k=\frac{\sum_i w_{i,k}r_{i,k}}{\sum_i w_{i,k}},\qquad \hat r_{i,k}=w_{i,k}r_{i,k}+(1-w_{i,k})\,m_k .
\]
Suffix returns $\hat R_{i,k}=\sum_{j\ge k}\gamma^{j-k}\hat r_{i,j}$ and raw returns $R_{i,k}$ (same
with $r$). Centred returns $C_{i,k}=\hat R_{i,k}-\frac1G\sum_{i'}\hat R_{i',k}$. ``Std'' is the
unbiased std (as in the code), so $\sum_i C_{i,k}^2=(G-1)\,\mathrm{Std}_i(\hat R_{i,k})^2$.

\begin{itemize}
\item \textbf{A} (eq.~\eqref{method:eq:adv}, default): $\hat A_{i,k}=C_{i,k}/(\hat s_k+\epsilon)$ with
  $\hat s_k=\mathrm{Std}_i\hat R_{i,k}$.
\item \textbf{B} (ablation): $\hat A_{i,k}=C_{i,k}/(s_k+\epsilon)$ with $s_k=\mathrm{Std}_i R_{i,k}$.
\end{itemize}

One identity is used throughout. By P\ref{method:prop:neutral}~(b), the group mean of
$\hat r_{\cdot,k}$ is $m_k$, so
\begin{equation*}
\hat r_{i,k}-m_k=w_{i,k}\,(r_{i,k}-m_k).\tag{$*$}\label{trust:eq:star}
\end{equation*}

\subsection{Test 5: Lemma 1 holds after imputation, for A and B}

Add $c_j$ to every reward of the group at chunk $j$. Then $m_j\to m_j+c_j$, so
$\hat r_{i,j}\to w(r+c_j)+(1-w)(m_j+c_j)=\hat r_{i,j}+c_j$. Every imputed reward moves by the same
$c_j$, as do the raw rewards. Every $\hat R_{i,k}$ and every $R_{i,k}$ then moves by the same
$\sum_{j\ge k}\gamma^{j-k}c_j$. So $C$, $\hat s_k$ and $s_k$ are unchanged, and so are the advantages
under A and under B.

The shift may differ between groups, because everything is computed per group.
\textbf{Condition:} the done boundaries must be shared within a group (fixed horizon). A member
whose episode ended before chunk $j$ would not receive $c_j$. The test uses dones only at the end;
whether production dones are always shared within a group is \chk.

\subsection{Test 6, option A: the total is fixed at \texorpdfstring{$G-1$}{G-1}}

$\sum_i \hat A_{i,k}^2=(G-1)\,\hat s_k^2/(\hat s_k+\epsilon)^2=(G-1)\,(1-O(\epsilon/\hat s_k))$,
whatever the weights.

Distrust therefore cannot change the total; it only redistributes it. By \eqref{trust:eq:star},
distrusted members' centred rewards shrink toward 0, so the trusted members carry a larger share.
If $G_T$ members are trusted and typical, each has $\hat A^2\approx (G-1)/G_T$ instead of
$(G-1)/G$: a factor $\sqrt{G/G_T}$, which is 2 for 12 of 16 distrusted.

\subsection{Test 6, option B: a fixed scale and a bound at the last chunk}

\textbf{(i) The scale does not depend on the weights.} $s_k$ is computed from the raw returns, so
$\hat A_{i,k}(s_k+\epsilon)=C_{i,k}$ exactly. The weights move the centre and the imputed values,
never the scale.

\textbf{(ii) At the last chunk, distrust never increases the total.} At $k=N-1$,
$C_{i,k}=\hat r_{i,k}-m_k=w_{i,k}(r_{i,k}-m_k)$ by \eqref{trust:eq:star}. Write $\bar r_k$ for the
plain mean. Then
\[
\sum_i C_{i,k}^2=\sum_i w_i^2(r_i-m)^2\ \le\ \sum_i w_i(r_i-m)^2\ \le\ \sum_i w_i(r_i-\bar r)^2\ \le\
\sum_i (r_i-\bar r)^2=(G-1)s_k^2 .
\]
The first and last steps use $w\le1$. The middle step holds because $m$ minimises
$\sum_i w_i(r_i-c)^2$ over $c$. Hence $\sum_i\hat A_{i,k}^2\le (G-1)s_k^2/(s_k+\epsilon)^2$, the
value at full trust. The same holds at every chunk when $\gamma=0$.

\textbf{(iii) At earlier chunks the bound fails.} $C_{i,k}=\sum_{j\ge k}\gamma^{j-k}w_{i,j}(r_{i,j}-m_j)$
(from \eqref{trust:eq:star}, using $\sum_i(\hat r_{i,j}-m_j)=0$). Each chunk's vector shrinks, but
cancellation across chunks in the raw returns can be undone.

Example: $G=2$, $r_0=(1,-1)$, $r_1=(-1,1)$, $\gamma=0.95$, $w_1=0.05$ for both members.
\begin{itemize}
\item Raw: $C_0=\pm0.05$.
\item Imputed: $C_0=\pm0.9525$.
\item The total grows about 360 times ($|\hat A|$ goes from 0.71 to 13.5).
\end{itemize}

\subsection{What does not hold: per-member monotonicity}

``The trusted members' advantages do not grow when others are distrusted'' is false under both
policies. Example: $G=16$, rewards eight 0 and eight 10, seven of the zeros distrusted with
$w=0.05$. The trusted zero-member's $|\hat A|$ is:

\begin{center}
\begin{tabular}{@{}lccc@{}}
\toprule
 & Full trust & B & A\\
\midrule
trusted zero-member $\lvert\hat A\rvert$ & 0.97 & 1.66 & 3.47\\
\bottomrule
\end{tabular}
\end{center}

Distrust moves $m_k$ toward the trusted majority, so away from a trusted minority member. B keeps
the scale, but the centre still moves. The test
\code{test\_d1\_counterexamples\_per\_member\_growth\_and\_earlier\_chunks} pins this example and the
one in (iii).

\subsection{Side result: exact neutrality of the advantages holds only under A}

With $w_{i,k}=0$, $r_{i,k}$ drops out of $\hat r$ entirely (P\ref{method:prop:neutral}~(c)), so out
of $C$ for both policies. Under A it then drops out of the advantages too. Under B it still enters
through $s_k$, so changing it rescales the whole group. With $w_{\min}>0$ neither policy is exact.

\subsection{Summary}

\begin{center}
\small
\begin{tabularx}{\linewidth}{@{}lYY@{}}
\toprule
 & A (eq.~\eqref{method:eq:adv}, default) & B (ablation)\\
\midrule
Scale & std of imputed returns, depends on $w$ & std of raw returns, independent of $w$\\
$\sum_i \hat A^2$ per chunk & $= G-1$ always & $\le$ full-trust value at the last chunk; can grow earlier\\
Effect of heavy distrust & same total update, concentrated on fewer members & smaller total update at the last chunk\\
$w=0$ reward & no effect on the advantages & no effect on $C$; still sets the scale\\
Per member & can grow (e.g.\ 3.6$\times$ above) & can grow (e.g.\ 1.7$\times$ above)\\
\bottomrule
\end{tabularx}
\end{center}

Interpretation for the optimisation notes \chk:
\begin{itemize}
\item Under A, the FPO surrogate keeps its overall magnitude whatever the trust. With little trust,
  the update is driven by a few members, so its variance rises.
\item Under B, heavy distrust shrinks the last-chunk advantages, which is more conservative, at the
  cost of a reward scale that still contains the distrusted rewards.
\end{itemize}
